% ============================== 5. АНАЛИЗ ==============================
\chapter{Анализ на експерименталните резултати}\label{ch:analysis}

\section{Отговори на изследователските въпроси}

\begin{keybox}[В1. Точност на трите алгоритъма за трите атаки]
И деветте модела постигат Accuracy над 99,5\%. За DDoS тя е \val{RF-DDoS-Accuracy} (RF), \val{SVM-DDoS-Accuracy} (SVM) и \val{XGB-DDoS-Accuracy} (XGBoost). За BruteForce е \val{RF-BruteForce-Accuracy}, \val{SVM-BruteForce-Accuracy} и \val{XGB-BruteForce-Accuracy}, а за Bot \val{RF-Bot-Accuracy}, \val{SVM-Bot-Accuracy} и \val{XGB-Bot-Accuracy}. Най-трудна е атаката Bot, а най-слаб е SVM. Двата ансамбъла от дървета са практически безгрешни при всички атаки. По учебния материал Accuracy е подходяща тук, защото тестът е балансиран~\cite{borovska2022}. Същото показват и F1, и Specificity.
\end{keybox}

\begin{keybox}[В2. Ефект на редукцията на размерността]
Двуетапната редукция свива вектора от \val{n-features} до \val{n-selected} признака (с 74\%), без да губи качество: разликите в F1 спрямо пълния вектор са в рамките на $\pm0{,}03$ пункта (E1). Подборът на признаци е поне толкова добър, колкото PCA, с по-малко измерения, и запазва интерпретируемите признаци (E5). Редукцията ускорява класификацията със SVM \val{SVM-DDoS-speedup} пъти при DDoS. Анализът A4 показва, че пет признака стигат за DDoS и BruteForce, а за Bot са нужни около 20.
\end{keybox}

\begin{keybox}[В3. Статистическа значимост и стабилност. Преобучение]
Разсейването по фолдове и bootstrap интервалите са много малки (полуширина на F1 под $0{,}05$ пункта). Тестът на McNemar показва значима разлика в \val{mcn-n-sig} от \val{mcn-n} двойки. Практически значима е само разликата на SVM при Bot: \val{mcn-Bot-RF-SVM-b} записа, които RF класифицира вярно, а SVM грешно, срещу \val{mcn-Bot-RF-SVM-c} обратни ($p={}$\val{mcn-Bot-RF-SVM-p}). Преобучение няма: разликата между точността при обучение и при тест е най-много \val{gap-max} пункта, а оценката от кръстосаната валидация съвпада с тестовата (\cref{tab:overfit}).
\end{keybox}

\begin{keybox}[В4. Качество при реално съотношение и избор на праг]
При естествения дял на атаките (\val{nat-share-Bot}--\val{nat-share-DDoS}) Precision остава между \val{SVM-Bot-nat-Precision} (SVM при Bot) и 100\%. Проблемът на базовата честота (\cref{eq:baserate}) не се проявява, защото делът на фалшивите аларми е от порядъка на $10^{-4}$--$10^{-3}$, много по-малък от дела на атаките. При RF и XGBoost Precision 99\% се постига с практически пълен Recall, така че изборът на праг не изисква компромис.
\end{keybox}

\begin{keybox}[В5. Носители на информацията и от какво зависи точността]
BruteForce се разпознава по структурата на кратките сесии за вход: \val{ks-top5names-BruteForce} (средно $D=\val{ks-top5-BruteForce}$). DDoS се разпознава по размера на отговорите на сървъра и по TCP прозореца: \val{ks-top5names-DDoS}. Bot се разпознава по размера на пакетите и по интервалите между тях: \val{ks-top5names-Bot}, а пермутационната важност извежда и флага RST. И при трите атаки признаците на отделния поток носят отпечатъка на атаката. При Port Scan в CICIoT2023 не е така, затова там точността остава около 85\%. Постижимата точност се определя преди всичко от данните, а не от алгоритъма.
\end{keybox}

\begin{keybox}[В6. Приложимост извън случайното разделяне]
Точността не идва само от отпечатъка на инструмента: без признаците на TCP стека тя спада най-много с \val{rob-fp-maxdrop} пункта (A10). При хронологично разделяне DDoS и BruteForce остават над \val{rob-chrono-min-DDoSBF}, но при Bot Random Forest и XGBoost пропускат около 10\% от по-късните атаки, докато SVM открива \val{rob-chronological-all-SVM-Bot-sens} (A11). Комбинираният детектор XGBoost или SVM възстановява откриването до \val{rob-chronological-all-XGBorSVM-Bot-sens} (A14). Моделите не се пренасят между инструменти с различен протокол (A12), а фалшивите аларми зависят силно от това дали нормалният трафик при обучението е разнообразен (A13).
\end{keybox}

\section{Защо трите атаки се откриват толкова точно}
Резултатите се обясняват с начина, по който всяка атака променя отделния мрежов поток. Опитът за вход при BruteForce е кратка сесия с еднаква структура: еднакъв брой пакети, еднакъв размер на заглавките и на полезните данни, еднакви TCP опции. Нормалният потребител рядко създава хиляди такива еднакви сесии. Затова признак като \texttt{Fwd Seg Size Min} сам разделя класовете. DDoS с HOIC изпраща HTTP заявки с типичен размер, а сървърът отговаря с типичен размер. Флагове и размер на TCP прозореца, характерни за инструмента, се срещат рядко в нормалния трафик. Botnet Ares поддържа периодична връзка с команден сървър: малки, еднакви по размер заявки и отговори през регулярни интервали, а броят на пакетите с флаг RST е сред най-важните признаци за моделите.

Сканирането на портове в CICIoT2023 е обратният случай. Отделната проба към един порт е нормална кратка TCP сесия. Аномалията се вижда само при агрегиране по източник, когато един адрес се свързва с много различни портове за кратко време. Признак от този тип липсва в записите. Това сравнение е основният методологичен извод на изследването: преди да се търси по-добър алгоритъм, трябва да се провери дали признаците на записа изобщо съдържат информация за атаката.

\section{Сравнение на алгоритмите}
\begin{table}[H]
\centering\small
\caption{Обобщено сравнение на Random Forest, SVM и XGBoost в изследването}\label{tab:algo-compare}
\begin{tabularx}{\textwidth}{@{}l L L L@{}}
\toprule
\textsf{\textbf{Критерий}} & \textsf{\textbf{Random Forest}} & \textsf{\textbf{SVM (RBF)}} & \textsf{\textbf{XGBoost}}\\
\midrule
Качество при DDoS и BruteForce & практически безгрешен & над 99,9\% & практически безгрешен\\
Качество при Bot & \val{RF-Bot-Accuracy} & \val{SVM-Bot-Accuracy} & \val{XGB-Bot-Accuracy}\\
Мащабируемост на обучението & почти линейна, паралелизира се & $O(n^2)$--$O(n^3)$, нужна е подизвадка & почти линейна, хистограмен алгоритъм\\
Скорост в експлоатация (Bot) & \val{RF-Bot-throughput} записа/s & \val{SVM-Bot-throughput} записа/s & \val{XGB-Bot-throughput} записа/s\\
Нужда от мащабиране & не & задължително (\cref{eq:zscore}) & не\\
Защита от преобучение & bagging, листа с мин. брой записи & марж, параметър $C$ & регуляризация $\Omega$, subsample\\
Обяснимост & MDI и пермутационна важност & само пермутационна важност & gain и пермутационна важност\\
\bottomrule
\end{tabularx}
\end{table}

За практическа IDS в център за данни XGBoost е най-добрият избор. Той е поне толкова точен, колкото Random Forest, класифицира около пет пъти по-бързо от него и от 2,6 до 17 пъти по-бързо от SVM. Освен това дава обяснима важност на признаците, която тук се използва и за подбора им. Random Forest е почти равностоен и по-прост за настройка. SVM е най-слабият избор при случайно разделяне: по-ниска точност при Bot, по-бавна класификация и необходимост от подизвадка при обучението. Проверките за приложимост обаче показват и силна страна на SVM: гладката граница го прави по-устойчив при промяна на данните. Той открива новите потоци на Bot, които дърветата пропускат (A11), и не вдига масово фалшиви аларми върху необичаен нормален трафик (A13). Затова за Bot препоръчваме комбинирания детектор XGBoost или SVM, когато се очакват нови разновидности на атаката. Част от изоставането на SVM при Bot се дължи на по-малката обучаваща извадка (\num{50000} срещу \val{train-rows-Bot} записа), но кривата на обучение показва, че и при равен брой записи SVM е под ансамблите (\texttt{results/learning\_curves.csv}).

\section{Редукция на размерността и практическа приложимост}
Подборът на признаци е по-подходящ от PCA за IDS по три причини. Първо, той отчита етикета и запазва признаците, които различават атаката, докато PCA запазва дисперсията, която може да е безполезна. Второ, запазените признаци са оригинални величини (размер на пакет, флаг, TCP прозорец), които операторът разбира. Трето, IDS трябва да изчислява само \val{n-selected} от \val{n-features} признака, което намалява натоварването на сензора. Филтърът по дисперсия е полезен винаги, защото CICFlowMeter-V3 записва някои признаци винаги с нула~\cite{guyon2003}.

Експериментът E2 показва, че при атаки, които са 14--28\% от трафика, и при много малък дял фалшиви аларми, моделите остават полезни и при естествено съотношение. В реален център за данни делът на атаките обикновено е много по-малък. Тогава дори FPR от порядъка на $10^{-4}$ може да даде повече фалшиви, отколкото истински аларми (\cref{eq:baserate}). Практически мерки са изборът на праг по бюджет от аларми (\cref{tab:ops}), агрегирането на алармите по източник във времеви прозорец, отделен модел за всяка категория атаки, тъй като A7 показва, че моделите не се пренасят между атаките, и периодично преобучение с нови данни~\cite{axelsson2000,sommer2010}.

\section{Заплахи за валидността и насоки за развитие}
\paragraph{Вътрешна валидност.} Дубликатите и противоречивите вектори са премахнати преди разделянето, а цялата обработка се обучава само върху обучаващите данни. Признаците \texttt{Timestamp} и \texttt{Dst Port} са изключени, за да не може моделът да разпознава атаката по момента или по порта. Основният тест използва случайно разделяне на записите от едни и същи дни. Затова са добавени хронологично разделяне (A11) и тест на нормален трафик от непознат ден (A13), които са по-близо до реалната експлоатация~\cite{arp2022}. Те потвърждават резултатите за DDoS и BruteForce, а за Bot показват спад при ансамблите, който се компенсира с комбинирания детектор.

\paragraph{Външна валидност и валидност на конструкта.} Атаките в CSE-CIC-IDS2018 са проведени с конкретни инструменти (HOIC, LOIC, Patator, Ares). Част от разделящите признаци, като началния размер на TCP прозореца и минималния размер на сегмента, отразяват мрежовия стек на атакуващите машини, а не самата атака. Друг инструмент или друга операционна система могат да променят тези стойности. Проверката A10 показва, че без тези признаци точността почти не се променя, но A12 показва, че модел, обучен върху един инструмент, не открива друг инструмент с различен протокол. Известни са и грешки в етикетите и в извличането на потоците в наборите на CIC~\cite{engelen2021,liu2022}. Тук те са ограничени с премахването на противоречивите вектори, но не могат да се изключат напълно. Затова резултатите трябва да се четат като горна граница за откриването на тези инструменти, а не като гаранция за всяка атака от същия тип.

\paragraph{Насоки за развитие.} Най-важните посоки са: тест върху атаки от същия тип, проведени в други дни или с други инструменти; тест върху друг реален набор (например CIC-IDS2017 или CICIoT2023) за оценка на преносимостта; многокласова класификация на всички 14 атаки в CSE-CIC-IDS2018; обучение на целия набор (около 16 милиона потока) с разпределени инструменти за големи данни~\cite{leevy2020}; и нови признаци, агрегирани по източник, които биха позволили откриването и на Port Scan~\cite{neto2023,sharafaldin2018}.
